\section*{5. 技术直觉}
Working-set 分析背后的想法并不复杂：一次昂贵的删除操作，必须能在图的顺序自由度中找到对应的信息量。一个顶点在堆中停留很久，删除它就贵；但停留很久也说明在它进入堆和离开堆之间有许多顶点参与竞争，这些顶点可能形成更多合法的距离顺序。高代价在这里有对应的信息来源，最终由 \(\log D\) 这样的参数承担。

Bottleneck 部分的直觉则相反：如果某些顶点由层次和支配关系几乎确定，算法就不该继续用普通堆比较来“发现”它们。Dijkstra with lookahead 把连续的 unmarked bottleneck 当作一段窄通道处理；recursive Dijkstra 则把 bottleneck 视为局部源点，遇到它就开启新的 Dijkstra 运行。两种写法不同，但都利用了同一个事实：低层 bottleneck 会支配更高层区域。

Heap 的实现只是把这种记账方式具体化：新元素留在较小层，旧元素逐步进入更大层，因此元素越旧，可用于支付删除代价的 working set 也越大。双指数大小阈值保证 inner heap 的数量很少，外层维护不会成为新的瓶颈。
